%% ch07-qec-security.tex — Chapter VII: Auditing Error Correction on Shared Cloud Hardware
%% Source: guo2026qecaudit (arXiv:2608.26600); artifact anonymous.4open.science usenix_qec_hse
%% Label prefix: qec:
%% STE edit (ASD-STE100 style, ~80% strictness) of chapters/ch07-qec-security.tex: prose shortened and
%% split; every number, citation, cross-reference, equation, environment, table and figure unchanged.
\chapter{Auditing Error Correction on Shared Cloud Hardware}\label{ch:qec}

This chapter addresses the first half of RQ5: how do we make error-corrected computation on shared cloud quantum processors auditable and secure? The dissertation's thesis is that a co-designed classical stack surrounds the QPU at every layer. One of those layers, this chapter argues, must be a \emph{security} layer attached at the error-correcting encoder. A fixed, publicly known encoder circuit is a reusable target for a co-located adversary who can inject a single Pauli fault. An encoder drawn afresh from a seeded ensemble at every run changes the map from physical faults to logical effects. The syndrome check then becomes a defense. We make this argument quantitative as follows:
\begin{itemize}
  \item define the \emph{accepted logical disturbance} (\cref{sec:qec:disturbance});
  \item derive its expectation over Haar-random encoders (\cref{sec:qec:haar});
  \item replace the exponentially costly Haar ensemble by a polynomial-cost seeded Clifford ensemble (\cref{sec:qec:hse});
  \item specify an audit protocol (\cref{sec:qec:protocol});
  \item evaluate it by exact enumeration up to eight qubits, with the fixed \code{5}{1}{3} code as control, and in gate-level simulations up to $n=17$ qubits (\cref{sec:qec:experiments}).
\end{itemize}
\Cref{sec:qec:discussion} states the limitations, \cref{sec:qec:relation} connects the chapter to \cref{ch:synthesis,ch:pqc}, and \cref{sec:qec:summary} summarizes.

\section{Introduction}\label{sec:qec:intro}

Quantum processors are consumed as a cloud service. A user submits a circuit, and a scheduler places it on a device that other users' circuits also occupy or have just occupied. The result comes back with no attestation of what else ran on the same chip. Multi-programming of a single QPU has been proposed as a throughput measure~\cite{das2019multiprogramming}. For NISQ workloads, the security consequences of co-tenancy have been studied. One job can direct crosstalk at another~\cite{ashsaki2020crosstalk}, and defenses that scan for malicious circuit patterns are needed~\cite{deshpande2022antivirus}.

Error-corrected workloads change the picture in two ways. First, the encoder that maps a logical state into a code is a long, structured circuit that runs identically every time. If it is public, an adversary can precompute which physical fault produces which logical effect. Second, error correction already contains a defense: syndrome extraction followed by correction or rejection catches exactly the low-weight disturbance a fault injector can produce. The question is how to stop the adversary from routing around it.

Fault-injection attacks are classical territory. Boneh, DeMillo and Lipton showed that a single induced fault can break RSA signatures~\cite{boneh1997faults}, and cryptographic-module standards have since required fault detection and randomized countermeasures~\cite{nist2019fips1403}. The analogue we propose for QEC is \emph{reseeding}: the encoder is drawn from a seeded ensemble at each run. A fault tuned to one encoder then maps to a detectable syndrome under the next.

The idea resembles randomized compiling~\cite{wallman2016randomizedcompiling}, which twirls gates to turn coherent noise into stochastic noise. It differs in intent and in analysis. Our adversary is a strategist, not noise. Our figure of merit is the logical damage that survives the syndrome check, not the average gate fidelity.

Reseeding is more than a slogan only if three questions are answered.
\begin{enumerate}
  \item What quantity does the adversary want to maximize and the defender want to bound?
  \item What does the ideal defense, a Haar-random encoder, achieve, and can a polynomial-size circuit achieve nearly the same?
  \item How does an operator \emph{audit} that the randomness is really there, in the vocabulary of the standards that govern randomness in cryptographic modules~\cite{nist2015sp80090a,nist2019fips1403}?
\end{enumerate}
The contributions answer them in turn.
\begin{enumerate}
  \item A bounded cloud fault model (\cref{sec:qec:threat}) and the accepted logical disturbance (\cref{sec:qec:disturbance}), an acceptance-weighted measure of the harmful logical action that survives syndrome-based rejection. Its Haar expectation has an exact expression that depends only on second moments, so it holds for any unitary 2-design (\cref{sec:qec:haar}).
  \item The Hadamard-structured ensemble (HSE), a seeded family of Clifford encoders built from Hadamard, phase, permutation and CNOT layers of depth $d$, with polynomial-cost derivation of stabilizers and logicals (\cref{sec:qec:hse}).
  \item An audit protocol (\cref{sec:qec:protocol}) and its evaluation (\cref{sec:qec:experiments}). For eight-qubit HSE encoders of depth $12$ with all $3n=24$ weight-one Pauli faults, reseeding reduces the mean accepted logical disturbance from $0.150$ (fault chosen after the encoder is known) to $0.020$ (fault committed before the seed is drawn). This $86.7\%$ reduction is attributable to rejection. Also, $18.5\%$ of the sampled encoders satisfy the exact error-correction conditions, the fixed \code{5}{1}{3} code corrects every tested fault, and gate-level Stim simulations at $n\in\{5,9,13,17\}$ confirm the scaling.
\end{enumerate}

\section{Preliminaries}\label{sec:qec:prelim}

\subsection{Stabilizer Codes and the Five-Qubit Code}\label{sec:qec:stabilizer}

Here we fix the notation of the stabilizer formalism reviewed in \cref{sec:bg:qec}. An $\code{n}{k}{\delta}$ stabilizer code is defined by an abelian subgroup $\mathcal{S}$ of the $n$-qubit Pauli group that does not contain $-I$ and has $n-k$ independent generators $g_1,\dots,g_{n-k}$. The code space is
\begin{equation}\label{eq:qec:codespace}
  \mathcal{C}=\{\ket{\phi}:\ g\ket{\phi}=\ket{\phi}\ \ \forall g\in\mathcal{S}\},\qquad \dim\mathcal{C}=2^{k}.
\end{equation}
A Pauli $P$ has syndrome $\sigma(P)\in\{0,1\}^{n-k}$ with $\sigma_j(P)=0$ if $P$ commutes with $g_j$ and $1$ otherwise. Paulis with trivial syndrome form the normalizer $N(\mathcal{S})$. Those in $N(\mathcal{S})\setminus\mathcal{S}$ act as non-trivial logical operators, and the distance $\delta$ is the minimum weight of such an operator. The Knill--Laflamme conditions~\cite{knill1997theory} state that a set of errors $\mathcal{F}$ is correctable if and only if
\begin{equation}\label{eq:qec:kl}
  \Pi\,P_a^{\dagger}P_b\,\Pi \;=\; c_{ab}\,\Pi\qquad\forall\,P_a,P_b\in\mathcal{F}\cup\{I\},
\end{equation}
with $\Pi$ the projector onto $\mathcal{C}$ and $c_{ab}$ scalars. For a stabilizer code and Pauli errors, \cref{eq:qec:kl} holds when every product $P_a^{\dagger}P_b$ either lies in $\mathcal{S}$ or anticommutes with some generator. The weaker condition that every $P\in\mathcal{F}$ be \emph{detectable} requires only that $P\notin N(\mathcal{S})\setminus\mathcal{S}$.

The \code{5}{1}{3} code~\cite{laflamme1996perfect,bennett1996mixed} is the smallest code that corrects an arbitrary single-qubit error. Its four generators are the cyclic shifts of $XZZXI$, and its logical operators can be taken as $\bar X=X^{\otimes 5}$ and $\bar Z=Z^{\otimes 5}$~\cite{nielsen2010qcqi} (\cref{tab:qec:code513}). The code is \emph{perfect}: its $2^{4}-1=15$ non-zero syndromes are in one-to-one correspondence with the $3n=15$ weight-one Pauli faults, so every such fault is detected and uniquely identified. This makes the code the right test bed for the audit. An adversary restricted to weight-one faults cannot defeat the ideal \code{5}{1}{3} encoder. Any surviving disturbance must therefore be attributed to the \emph{encoder ensemble}, not to the code.

\begin{table}[htbp]
\caption{Stabilizer generators and logical operators of the \code{5}{1}{3} code (textbook form, after \cite{laflamme1996perfect,nielsen2010qcqi}). Each of the $15$ weight-one Pauli faults has a distinct non-zero syndrome.}
\label{tab:qec:code513}
\centering\small
\begin{tabular}{@{}lccccc l@{}}
\toprule
Operator & $q_1$ & $q_2$ & $q_3$ & $q_4$ & $q_5$ & Role \\
\midrule
$g_1$ & $X$ & $Z$ & $Z$ & $X$ & $I$ & stabilizer generator \\
$g_2$ & $I$ & $X$ & $Z$ & $Z$ & $X$ & stabilizer generator \\
$g_3$ & $X$ & $I$ & $X$ & $Z$ & $Z$ & stabilizer generator \\
$g_4$ & $Z$ & $X$ & $I$ & $X$ & $Z$ & stabilizer generator \\
\midrule
$\bar X$ & $X$ & $X$ & $X$ & $X$ & $X$ & logical $X$ (weight 5; weight-3 representatives exist) \\
$\bar Z$ & $Z$ & $Z$ & $Z$ & $Z$ & $Z$ & logical $Z$ (weight 5; weight-3 representatives exist) \\
\midrule
\multicolumn{7}{@{}l}{Weight-one faults: $\{X_i,Y_i,Z_i\}_{i=1}^{5}$, $15$ in total; syndromes $\sigma\in\{0,1\}^{4}\setminus\{0000\}$, $15$ in total.} \\
\bottomrule
\end{tabular}
\end{table}

\subsection{Encoders, Syndromes and Rejection}\label{sec:qec:encoders}

An \emph{encoder} for a $k=1$ code is a unitary $E$ on $n$ qubits that maps the logical qubit $A$, together with $n-1$ ancillas $B$ prepared in $\ket{0}$, into the code space,
\begin{equation}\label{eq:qec:encoder}
  E\big(\ket{\psi}_A\otimes\ket{0^{\,n-1}}_B\big)\in\mathcal{C},\qquad
  \mathcal{S}=\big\{E\,Z_j\,E^{\dagger}:\ j=2,\dots,n\big\},\qquad
  \bar X=E X_1 E^{\dagger},\ \ \bar Z=E Z_1 E^{\dagger}.
\end{equation}
When $E$ is Clifford, the stabilizers and logicals in \cref{eq:qec:encoder} are Paulis. Tableau simulation~\cite{aaronson2004stabilizer} computes them in time polynomial in $n$ and in the gate count of $E$. The code projector is $\Pi_E=E\,\Pi_0\,E^{\dagger}$ with $\Pi_0=I_A\otimes\ket{0^{\,n-1}}\!\bra{0^{\,n-1}}_B$. Up to the choice of generators, the syndrome measurement is the measurement of $\{Z_j\}_{j\ge 2}$ in the decoded frame.

We consider \emph{syndrome-based rejection}: a run whose syndrome is non-trivial is discarded. Rejection is the right audit primitive for two reasons. The accepted output depends on the fault only through $\Pi_E$. And an operator who rejects rather than corrects cannot be misled by a decoder that the adversary has also studied. Correction is the operational goal, and \cref{sec:qec:discussion} discusses how the audit extends to it.

\subsection{Random Unitaries, Designs and the Clifford Group}\label{sec:qec:designs}

A Haar-random encoder is the natural idealization of ``the adversary knows nothing about the code.'' Haar-random unitaries on $n$ qubits cannot be implemented with polynomially many gates. But unitary $t$-designs, ensembles whose first $t$ moments agree with Haar's, reproduce many of their statistical properties~\cite{dankert2009designs,gross2007designs}.

The uniform distribution over the $n$-qubit Clifford group is a unitary 2-design and indeed a 3-design~\cite{webb2016clifford3design}. A uniformly random Clifford element can be sampled and compiled in polynomial time~\cite{koenig2014clifford,bravyi2021hadamardfree}. \Cref{sec:qec:haar} exploits one consequence: any quantity quadratic in the encoder and its adjoint has the same expectation under Haar and under any 2-design. The accepted logical disturbance is such a quantity.

\section{The Bounded Cloud Fault Model}\label{sec:qec:threat}

\Cref{fig:qec:threat} depicts the setting. A client prepares a logical state, chooses (or is assigned) an encoder, and submits the encoded computation to a shared QPU. A co-tenant adversary occupies neighbouring qubits during a \emph{fault window} between encoding and syndrome extraction. The syndrome is then measured, and the run is accepted or rejected. The adversary is bounded as follows.

% alt: Threat-model diagram. On the left a client generates a seed from a validated random-bit generator and derives a seeded HSE encoder. Inside a dashed box labelled shared cloud QPU, the encoder is applied, a fault window follows in which a co-tenant adversary injects a weight-one Pauli fault, then syndrome extraction leads to an accept-or-reject decision; accepted runs produce the logical output, rejected runs are discarded. A note contrasts a fixed public encoder, which the adversary learns once, with per-run reseeding.
\begin{figure}[htbp]
\centering
\begin{tikzpicture}[
  font=\small,
  box/.style={draw,rounded corners=2pt,align=center,minimum height=10mm,text width=27mm,inner sep=3pt,fill=white},
  trusted/.style={box,fill=cbSky!20,draw=cbBlue},
  adv/.style={box,fill=cbRed!15,draw=cbRed},
  dec/.style={draw,diamond,aspect=2.2,align=center,inner sep=1pt,fill=white,font=\footnotesize},
  arr/.style={-{Latex[length=2.2mm]},thick}]
  \node[trusted] (seed) {Seed $s$\\validated DRBG};
  \node[trusted,right=6mm of seed] (enc) {HSE encoder\\$E_s=\textsc{HSE}(s,d)$};
  \node[box,right=12mm of enc] (encode) {Encode\\$E_s(\ket{\psi}\ket{0^{n-1}})$};
  \node[box,right=6mm of encode] (window) {Fault window\\$P\in\mathcal{F}$};
  \node[adv,above=7mm of window] (advn) {Co-tenant adversary\\weight-one Pauli};
  \node[box,below=9mm of window] (synd) {Syndrome\\extraction $\sigma$};
  \node[dec,left=8mm of synd] (dec) {$\sigma=0$?};
  \node[trusted,left=20mm of dec] (out) {Accept:\\logical output};
  \node[box,below=7mm of dec,text width=22mm,minimum height=8mm] (rej) {Reject: discard};
  \draw[arr] (seed) -- (enc);
  \draw[arr] (enc) -- (encode);
  \draw[arr] (encode) -- (window);
  \draw[arr,cbRed] (advn) -- (window);
  \draw[arr] (window) -- (synd);
  \draw[arr] (synd) -- (dec);
  \draw[arr] (dec) -- node[pos=0.62,above,font=\footnotesize]{yes} (out);
  \draw[arr] (dec) -- node[right,font=\footnotesize]{no} (rej);
  \draw[arr,dashed,cbRed] (advn.west) -| node[pos=0.42,above,font=\footnotesize]{learns $E$ if fixed} (enc.north);
  \begin{scope}[on background layer]
    \node[draw=cbGray,dashed,rounded corners=3pt,fit=(encode)(window)(synd)(dec)(rej)(advn),inner sep=5pt,label={[cbGray,font=\footnotesize]below:shared cloud QPU}] {};
  \end{scope}
\end{tikzpicture}
\caption{The bounded cloud fault model. The client derives a seeded encoder from a validated random-bit generator; on the shared QPU a co-tenant adversary injects one weight-one Pauli fault during the window between encoding and syndrome extraction; runs with a non-trivial syndrome are rejected. A fixed public encoder is learned once by the adversary (dashed arrow); per-run reseeding removes that knowledge. Schematic.}
\label{fig:qec:threat}
\end{figure}

\begin{assumption}[Bounded cloud fault model]\label{ass:qec:model}
\begin{enumerate}
  \item \emph{Fault capability.} During the fault window of one run the adversary can apply at most one weight-one Pauli $P\in\mathcal{F}=\{X_i,Y_i,Z_i\}_{i=1}^{n}$ to the encoded register, for instance through directed crosstalk or a mistimed control pulse~\cite{ashsaki2020crosstalk}. Native device noise is either absent (exact enumeration) or modelled as depolarizing noise on every gate (gate-level simulation).
  \item \emph{Knowledge.} The adversary knows the code family, the encoder ensemble and the audit protocol. In the \emph{adaptive} setting the adversary also learns the encoder $E$ of the current run before choosing $P$; this models a fixed, public encoder, since a fixed encoder is learned once and known thereafter. In the \emph{committed} setting the adversary must choose $P$ before the seed of the current run is drawn.
  \item \emph{Trusted components.} Seed generation, encoder derivation from the seed, syndrome extraction and the accept/reject decision are trusted and, in the enumeration experiments, ideal.
  \item \emph{Goal.} The adversary seeks to maximize the harmful logical action that survives rejection; the defender seeks to bound it.
\end{enumerate}
\end{assumption}

The model is deliberately bounded. A single weight-one fault is the smallest disturbance that a code is built to handle. It therefore isolates whether the \emph{encoder}, rather than the code, is the weak point. The adaptive-versus-committed distinction is the formal content of ``reseeding.'' \Cref{fig:qec:adversaries} shows the two settings: they differ only in whether the adversary chooses after or before the seed.

% alt: Two rows of four boxes. Top, the adaptive adversary: draw the seed and derive the encoder, the adversary learns the encoder and picks the worst fault, the syndrome check, and the outcome. Bottom, the committed adversary: the adversary commits to a fault first, then the seed is drawn and the encoder derived, the syndrome check, and the outcome. A note says that only the order of the first two steps differs.
\begin{figure}[htbp]
\centering
\begin{tikzpicture}[
  font=\footnotesize,
  box/.style={draw,rounded corners=2pt,align=center,minimum height=9mm,text width=27mm,inner sep=2pt},
  tr/.style={box,fill=cbSky!20,draw=cbBlue},
  ad/.style={box,fill=cbRed!12,draw=cbRed},
  res/.style={box,fill=cbOrange!15,draw=cbOrange},
  arr/.style={-{Latex[length=2mm]},thick}]
  \node[font=\footnotesize\bfseries,anchor=west,text=cbBlue] at (-1.5,0.9) {Adaptive (seed-aware): fixed or public encoder};
  \node[tr]  (a1) at (0,0)   {Draw seed, derive $E_s$};
  \node[ad]  (a2) at (3.6,0) {Adversary learns $E_s$, picks worst $P$};
  \node[box] (a3) at (7.2,0) {Syndrome check};
  \node[res] (a4) at (10.8,0){$\Delta=1$ if any fault escapes};
  \draw[arr] (a1)--(a2); \draw[arr] (a2)--(a3); \draw[arr] (a3)--(a4);
  \node[font=\footnotesize\bfseries,anchor=west,text=cbBlue] at (-1.5,-1.6) {Committed (seed-blind): reseeded encoder};
  \node[ad]  (b1) at (0,-2.5)   {Adversary commits to $P^\star$};
  \node[tr]  (b2) at (3.6,-2.5) {Draw seed, derive $E_s$};
  \node[box] (b3) at (7.2,-2.5) {Syndrome check};
  \node[res] (b4) at (10.8,-2.5){$\Delta=1$ only if $P^\star$ escapes};
  \draw[arr] (b1)--(b2); \draw[arr] (b2)--(b3); \draw[arr] (b3)--(b4);
  \node[font=\scriptsize,text=cbGray] at (5.4,-3.5) {only the order of the first two steps differs};
\end{tikzpicture}
\caption{The two adversaries of \cref{ass:qec:model}. The adaptive adversary chooses the fault after it learns the encoder, which models a fixed public encoder; the committed adversary chooses before the seed is drawn, which models per-run reseeding. Schematic.}
\label{fig:qec:adversaries}
\end{figure}

\section{Accepted Logical Disturbance}\label{sec:qec:disturbance}

\subsection{Definition}\label{sec:qec:definition}

We need one number for the harm a fault does \emph{after} rejection has had its chance. Two quantities matter: the probability that the faulty run is accepted, and the fidelity of the accepted logical state with the intended one. A fault that is always rejected or acts trivially on the logical qubit does no harm. A fault that is accepted and flips the logical qubit does maximal harm. We evaluate the measure on the logical qubit entangled with a reference qubit $R$. This makes it independent of the logical state and yields the entanglement fidelity of the induced logical channel.

\begin{definition}[Accepted logical disturbance]\label{def:qec:disturbance}
Let $E$ be an encoder for a $k=1$ code with code projector $\Pi_E=E\Pi_0E^{\dagger}$, let $\ket{\Phi^{+}}_{AR}=(\ket{00}+\ket{11})/\sqrt2$, let $\ket{\Phi_E}=(E\otimes I_R)\big(\ket{\Phi^{+}}_{AR}\otimes\ket{0^{\,n-1}}_B\big)$ be the encoded half of a maximally entangled pair, and let $P$ be a fault operator. Write $p_{\mathrm{acc}}(E,P)=\norm{(\Pi_E P\otimes I_R)\ket{\Phi_E}}^{2}$ for the acceptance probability and $F_e(E,P)=\abs{\bra{\Phi_E}(P\otimes I_R)\ket{\Phi_E}}^{2}/p_{\mathrm{acc}}(E,P)$ for the entanglement fidelity of the accepted state. The \emph{accepted logical disturbance} is
\begin{equation}\label{eq:qec:disturbance}
  \Delta(E,P)\;=\;p_{\mathrm{acc}}(E,P)\,\big(1-F_e(E,P)\big)
  \;=\;\norm{(\Pi_E P\otimes I_R)\ket{\Phi_E}}^{2}-\abs{\bra{\Phi_E}(P\otimes I_R)\ket{\Phi_E}}^{2}.
\end{equation}
\end{definition}

The second form in \cref{eq:qec:disturbance} follows from $\Pi_E\ket{\Phi_E}=\ket{\Phi_E}$ and is the one we compute. It is the probability that the fault is accepted \emph{and} the accepted state is orthogonal, within the code space, to the intended one. Several properties are immediate. $\Delta\in[0,1]$. $\Delta=0$ whenever $P$ is rejected with certainty or acts as a logical identity. $\Delta$ is acceptance-weighted, so a fault that slips through with probability $p_{\mathrm{acc}}$ and then does full damage contributes $p_{\mathrm{acc}}$, not one.

For an ensemble of encoders $\mathcal{E}$ and the two adversaries of \cref{ass:qec:model} (\cref{fig:qec:adversaries}) the quantities audited are
\begin{equation}\label{eq:qec:adversaries}
  \bar\Delta_{\mathrm{adapt}}(\mathcal{E})=\mathbb{E}_{E\sim\mathcal{E}}\Big[\max_{P\in\mathcal{F}}\Delta(E,P)\Big],\qquad
  \bar\Delta_{\mathrm{commit}}(\mathcal{E})=\max_{P\in\mathcal{F}}\ \mathbb{E}_{E\sim\mathcal{E}}\big[\Delta(E,P)\big],
\end{equation}
and the gap between them is the value of reseeding. For a single fixed encoder the two coincide: a fixed encoder offers no protection beyond the code itself. In the experiments of \cref{sec:qec:experiments} the committed adversary's fault is also averaged over $\mathcal{F}$, to report the typical case alongside the worst case. The value $0.020$ quoted there is the mean over faults and seeds, as reported in \cite{guo2026qecaudit}.

\subsection{Clifford Encoders and Pauli Faults}\label{sec:qec:cliffordcase}

For the encoders and faults used in the experiments, \cref{eq:qec:disturbance} collapses to an indicator.

\begin{lemma}\label{lem:qec:indicator}
Let $E$ be a Clifford encoder and $P\neq I$ a Pauli fault, and write $E^{\dagger}PE=\pm\,L\otimes Q$ with $L$ a Pauli on the logical qubit $A$ and $Q$ a Pauli on the ancillas $B$. Then $P$ is accepted if and only if $Q\in\{I,Z\}^{\otimes(n-1)}$, and
\begin{equation}\label{eq:qec:indicator}
  \Delta(E,P)=\mathbb{1}\big[Q\in\{I,Z\}^{\otimes(n-1)}\big]\cdot\mathbb{1}\big[L\neq I\big]\in\{0,1\}.
\end{equation}
\end{lemma}
\begin{proof}
In the decoded frame the faulty state is $(L\otimes Q\otimes I_R)(\ket{\Phi^{+}}\otimes\ket{0^{\,n-1}})$ up to a phase. If $Q$ contains an $X$ or $Y$ factor it maps $\ket{0^{\,n-1}}$ to an orthogonal computational basis state, so $\Pi_0$ annihilates the state and $p_{\mathrm{acc}}=0$. Otherwise $Q\ket{0^{\,n-1}}=\ket{0^{\,n-1}}$, so $p_{\mathrm{acc}}=1$ and $F_e=\abs{\bra{\Phi^{+}}(L\otimes I)\ket{\Phi^{+}}}^{2}=\abs{\Tr L/2}^{2}$, which equals $1$ if $L=I$ and $0$ otherwise.
\end{proof}

\Cref{fig:qec:indicator} draws the decision of \cref{lem:qec:indicator}, which we read in two ways. First, $\Delta(E,P)=1$ exactly when $P\in N(\mathcal{S}_E)\setminus\mathcal{S}_E$, that is, when $P$ is an undetectable logical of the code generated by $E$. Hence $\max_{P\in\mathcal{F}}\Delta(E,P)=1$ if and only if the code has distance one, and $\bar\Delta_{\mathrm{adapt}}$ is the fraction of distance-one draws. Second, $\mathbb{E}_E[\Delta(E,P)]$ is the fraction of encoders for which the particular fault $P$ is an undetectable logical, the most a committed adversary can hope for.

% alt: Decision flow for a Clifford encoder and a Pauli fault. The fault is conjugated into the decoded frame, which gives a logical part L and an ancilla part Q. If Q contains an X or Y factor, the syndrome is non-trivial, the run is rejected and the disturbance is zero. Otherwise the run is accepted: if L is the identity the fault is a stabilizer, harmless, and the disturbance is zero; if L is not the identity the fault is an undetectable logical and the disturbance is one.
\begin{figure}[htbp]
\centering
\begin{tikzpicture}[
  font=\small,
  box/.style={draw,rounded corners=2pt,align=center,minimum height=9mm,text width=40mm,inner sep=3pt,fill=white},
  fault/.style={box,draw=cbRed},
  safe/.style={box,fill=cbGreen!15,draw=cbGreen,text width=46mm},
  harm/.style={box,fill=cbRed!15,draw=cbRed,text width=46mm},
  dec/.style={draw=cbBlue,diamond,aspect=2.4,align=center,inner sep=1pt,fill=cbSky!20,font=\small},
  arr/.style={-{Latex[length=2.2mm]},thick},
  node distance=6mm]
  \node[fault] (p) {Weight-one Pauli fault $P$};
  \node[box,below=of p] (conj) {Decoded frame:\\$E^{\dagger}PE=\pm\,L\otimes Q$};
  \node[dec,below=of conj] (q) {$Q\in\{I,Z\}^{\otimes(n-1)}$?};
  \node[dec,below=8mm of q] (l) {$L\neq I$?};
  \node[harm,below=of l] (bad) {Accepted and harmful:\\$P\in N(\mathcal{S}_E)\setminus\mathcal{S}_E$, $\Delta(E,P)=1$};
  \node[safe,right=10mm of q] (rej) {Rejected ($\sigma\neq 0$):\\$p_{\mathrm{acc}}=0$, $\Delta(E,P)=0$};
  \node[safe,right=10mm of l] (harmless) {Accepted and harmless:\\$P\in\mathcal{S}_E$, $\Delta(E,P)=0$};
  \draw[arr] (p) -- (conj);
  \draw[arr] (conj) -- (q);
  \draw[arr] (q) -- node[right,font=\footnotesize]{yes} (l);
  \draw[arr] (q) -- node[above,font=\footnotesize]{no} (rej);
  \draw[arr] (l) -- node[right,font=\footnotesize]{yes} (bad);
  \draw[arr] (l) -- node[above,font=\footnotesize]{no} (harmless);
\end{tikzpicture}
\caption{Decision flow of \cref{lem:qec:indicator} for a Clifford encoder $E$ and a Pauli fault $P$. An $X$ or $Y$ factor in the ancilla part $Q$ gives a non-trivial syndrome and rejection. Otherwise the run is accepted, and the fault is harmful exactly when its logical part $L$ is not the identity. The adaptive adversary wins when any weight-one fault reaches the bottom box; the committed adversary wins only when its chosen fault does.}
\label{fig:qec:indicator}
\end{figure}

\subsection{Expectation over Haar-Random Encoders}\label{sec:qec:haar}

The Haar-random encoder is the ideal against which structured ensembles are measured. Its expected disturbance has a closed form.

\begin{proposition}[Haar expectation of the accepted logical disturbance]\label{prop:qec:haar}
Let $d=2^{n}$, $k=1$, and let $P$ be any fixed $n$-qubit operator with $\Tr P=0$ and $P^{2}=I$ (in particular any non-identity Pauli). Then
\begin{equation}\label{eq:qec:haar}
  \mathbb{E}_{E\sim\mathrm{Haar}}\big[\Delta(E,P)\big]\;=\;\frac{3d}{2(d^{2}-1)}\;=\;\frac{3\cdot 2^{n}}{2\,(4^{n}-1)},
\end{equation}
independently of $P$. The same value is obtained for any unitary 2-design, in particular for the uniform distribution over the $n$-qubit Clifford group. For $k$ logical qubits the constant becomes $d\,(2^{k}-2^{-k})/(d^{2}-1)$.
\end{proposition}
\begin{proof}[Proof sketch]
Set $V=E^{\dagger}PE$. Both terms of \cref{eq:qec:disturbance} are quadratic in $V$, hence in $E$ and $E^{\dagger}$, so only the second moment of $E$ enters and any 2-design gives the Haar value. For Haar $E$ and $\Tr P=0$, $\Tr P^{2}=d$, the standard second-moment identity gives $\mathbb{E}[V\otimes V]=(d\,\mathrm{SWAP}-I)/(d^{2}-1)$, and therefore $\mathbb{E}[VAV]=(d\Tr(A)\,I-A)/(d^{2}-1)$ for any operator $A$. The acceptance term is $\bra{\Phi}(V\Pi_0V\otimes I_R)\ket{\Phi}$ with $\ket{\Phi}=\ket{\Phi^{+}}\otimes\ket{0^{\,n-1}}$; taking $A=\Pi_0$, $\Tr\Pi_0=2$, gives $\mathbb{E}[p_{\mathrm{acc}}]=(2d-1)/(d^{2}-1)$. The fidelity term is $\mathbb{E}\abs{\Tr(V\rho)}^{2}$ with $\rho=\Tr_R\ket{\Phi}\!\bra{\Phi}=\tfrac12 I_A\otimes\ket{0^{\,n-1}}\!\bra{0^{\,n-1}}$, and $\mathbb{E}\abs{\Tr V\rho}^{2}=\Tr[(\rho\otimes\rho)\,\mathbb{E}(V\otimes V)]=(d\Tr\rho^{2}-(\Tr\rho)^{2})/(d^{2}-1)=(d/2-1)/(d^{2}-1)$. Subtracting yields $(2d-1-d/2+1)/(d^{2}-1)=3d/(2(d^{2}-1))$. For general $k$ replace $\Tr\Pi_0=2$ by $2^{k}$ and $\Tr\rho^{2}=1/2$ by $2^{-k}$. The full derivation is given in \cite{guo2026qecaudit}, whose Theorem~1 states the same values in the form $\mathbb{E}[p_{\mathrm{acc}}]=(KN-1)/(N^{2}-1)$ and $\mathbb{E}[c]=N(K^{2}-1)/(K(N^{2}-1))$ with $N=2^{n}$ and $K=2^{k}$.
\end{proof}

For $n=5$, \cref{eq:qec:haar} gives $96/2046\approx 0.0469$, for $n=8$ it gives $768/131070\approx 0.0059$, and for $n=17$ it gives $\approx 1.1\times 10^{-5}$. Asymptotically the Haar value decays as $\tfrac{3}{2}\,2^{-n}$. As a consistency check, under the uniform Clifford group $E^{\dagger}PE$ is a uniformly random non-identity Pauli. By \cref{lem:qec:indicator} the disturbance is one exactly when its ancilla part is $Z$-type ($2^{n-1}$ choices) and its logical part is non-identity ($3$ choices). This gives $3\cdot 2^{n-1}/(4^{n}-1)=48/1023$ for $n=5$, in agreement with \cref{eq:qec:haar}.

The Haar value is the committed-adversary disturbance of an ensemble with no structure at all. A structured ensemble may lie above it (if its structure leaks) or below it (if its structure happens to favour good codes). And the adaptive-adversary disturbance of any ensemble is at least its committed-adversary disturbance, because a maximum over faults dominates each fault's mean.

\section{The Hadamard-Structured Ensemble}\label{sec:qec:hse}

Haar-random encoders are unimplementable. Even uniformly random Cliffords, though polynomial to sample~\cite{koenig2014clifford}, produce circuits whose depth and connectivity are not under the operator's control. The HSE is a seeded Clifford ensemble designed for the audit: shallow, of fixed layer structure, and cheap to derive from a seed. It is also rich enough to leave the Hadamard-free subgroup of the Clifford group~\cite{bravyi2021hadamardfree}, whose elements map $Z$-type Paulis to $Z$-type Paulis. A Hadamard-free encoder therefore produces only $Z$-type stabilizers and cannot detect a single phase-flip fault.

Let $s$ be a seed and $\mathrm{DRBG}(s)$ a deterministic random-bit generator instantiated from it~\cite{nist2015sp80090a}. A depth-$d$ HSE encoder is
\begin{equation}\label{eq:qec:hse}
  E_s=\prod_{\ell=1}^{d}\ \mathrm{C}_{M_\ell}\;\Pi_{\sigma_\ell}\;\mathrm{S}_{B_\ell}\;\mathrm{H}_{A_\ell},
  \qquad
  \mathrm{H}_{A}=\bigotimes_{i\in A}H_i,\quad
  \mathrm{S}_{B}=\bigotimes_{i\in B}S_i,\quad
  \mathrm{C}_{M}=\prod_{(i\to j)\in M}\mathrm{CNOT}_{i\to j},
\end{equation}
where in layer $\ell$ the Hadamard subset $A_\ell\subseteq[n]$, the phase subset $B_\ell\subseteq[n]$, the qubit permutation $\sigma_\ell\in S_n$ and the CNOT matching $M_\ell$ (a set of disjoint ordered pairs) are all read from $\mathrm{DRBG}(s)$. \Cref{fig:qec:hse} shows one layer on five qubits. Every factor in \cref{eq:qec:hse} is Clifford, so the stabilizer generators and logical operators follow from \cref{eq:qec:encoder} by tableau conjugation at cost $\bigO{n}$ per gate~\cite{aaronson2004stabilizer}. That is $\bigO{d\,n^{2}}$ for the whole encoder, against $\bigO{4^{n}}$ to write down a Haar-random unitary. The permutation layer is free on hardware with all-to-all connectivity and costs a SWAP network otherwise. On a fixed coupling graph it can be restricted to permutations realizable in constant depth, the topology-compliance constraint that \cref{ch:synthesis} enforced through a reward term.

% alt: Quantum circuit on five wires showing one layer of the Hadamard-structured ensemble: a Hadamard layer acting on a random subset of qubits, a phase (S) layer on another subset, a permutation drawn as a swap of two wires, and a CNOT layer on a random matching of disjoint pairs. A box marked layers two to d follows, then a box marked syndrome extraction and accept or reject.
\begin{figure}[htbp]
\centering
\begin{quantikz}[row sep={7mm,between origins},column sep=5mm,font=\small]
\lstick{$\ket{\psi}$} & \gategroup[5,steps=1,style={dashed,rounded corners,fill=cbSky!12,inner xsep=2pt},background,label style={label position=above,anchor=south,yshift=0.15cm}]{$\mathrm{H}_{A_\ell}$} & \gate{S}\gategroup[5,steps=1,style={dashed,rounded corners,fill=cbOrange!12,inner xsep=2pt},background,label style={label position=below,anchor=north,yshift=-0.15cm}]{$\mathrm{S}_{B_\ell}$} & \gategroup[5,steps=1,style={dashed,rounded corners,fill=cbGreen!12,inner xsep=2pt},background,label style={label position=above,anchor=south,yshift=0.15cm}]{$\Pi_{\sigma_\ell}$} & \ctrl{1}\gategroup[5,steps=1,style={dashed,rounded corners,fill=cbPurple!12,inner xsep=2pt},background,label style={label position=below,anchor=north,yshift=-0.15cm}]{$\mathrm{C}_{M_\ell}$} & \gate[5,style={fill=white}]{\text{layers }2,\dots,d} & \gate[5,style={fill=white}]{\text{syndrome; reject if }\sigma\neq 0} & \\
\lstick{$\ket{0}$} & \gate{H} &          & \swap{1}  & \targ{}  &  &  & \\
\lstick{$\ket{0}$} &          & \gate{S} & \targX{}  & \ctrl{1} &  &  & \\
\lstick{$\ket{0}$} & \gate{H} &          &           & \targ{}  &  &  & \\
\lstick{$\ket{0}$} &          &          &           &          &  &  &
\end{quantikz}
\caption{One layer of a Hadamard-structured-ensemble encoder on $n=5$ qubits: a Hadamard layer on the seeded subset $A_\ell$, a phase layer on $B_\ell$, a qubit permutation $\sigma_\ell$ (drawn as a swap) and a CNOT layer on a seeded matching $M_\ell$, followed by the remaining $d-1$ layers and by syndrome extraction with rejection. Schematic; the subsets shown are one example draw.}
\label{fig:qec:hse}
\end{figure}

The seed budget of the ensemble is small and explicit. Each layer consumes:
\begin{itemize}
  \item $n$ bits for the Hadamard subset;
  \item $n$ bits for the phase subset;
  \item about $\log_2 n!\approx n\log_2 n$ bits for the permutation;
  \item at most $n\log_2 n$ bits for the matching.
\end{itemize}
A depth-$d$ encoder is therefore determined by $\bigO{d\,n\log n}$ output bits of the generator: a few hundred bits for the five-qubit experiments and a few thousand at $n=17$. This fits in a single request to an SP 800-90A generator, and it fixes the requirement on the seed source precisely. Reseeding protects against a committed adversary only if the adversary cannot predict those bits. A validated DRBG is certified to provide exactly that, and no circuit-level cleverness can substitute for it. Because the same seed always yields the same circuit, an auditor who holds the seed log can regenerate every encoder and recompute every reported syndrome.

The budget also bounds the entropy of the ensemble from above: at most $\bigO{d\,n\log n}$ bits, against the $\Theta(n^{2})$ bits needed to specify a uniformly random Clifford~\cite{koenig2014clifford}. This is why a shallow HSE is not a 2-design, and why its disturbance can differ from the value of \cref{prop:qec:haar} in either direction.

Not every HSE encoder generates a good code. A draw whose CNOT layers never touch some ancilla leaves that ancilla's $Z$ stabilizer of weight one. A draw with too few Hadamards leaves the logical operators concentrated on few qubits. The audit does not filter such draws away. It measures how often they occur and how much disturbance they admit, because an operator who reseeds at every run will also draw them.

As reported in \cite{guo2026qecaudit}, $18.5\%$ of the sampled depth-$12$ HSE encoders on eight qubits satisfy the exact Knill--Laflamme conditions \cref{eq:qec:kl} for all $24$ weight-one faults. The fraction for uniformly random Clifford encoders is $14.5\%$. Such draws are as capable against this fault set as a designed distance-three code such as the \code{5}{1}{3} code. A much larger fraction detect all such faults, which is what rejection needs. The depth $d$ trades cost against quality. Deeper ensembles approach the uniform Clifford distribution and hence the 2-design value of \cref{eq:qec:haar}. Shallow ensembles are cheaper and, as the next sections show, can beat that value against a committed adversary because their structure favours spreading.

\section{The Audit Protocol}\label{sec:qec:protocol}

An audit is a procedure that an operator, a client or a third party can run. It certifies that the deployed encoder ensemble delivers the disturbance bound it claims, and that the randomness feeding it has the quality the claim assumes. \Cref{alg:qec:audit} states the procedure, and \cref{tab:qec:protocol} maps each step onto the standards that govern it. The protocol has two modes, which mirror \cref{ass:qec:model} and differ only in when the adversary chooses the fault. It has two evaluation back-ends. Exact enumeration by tableau simulation serves Clifford encoders and Pauli faults, and sampled gate-level simulation with Stim~\cite{gidney2021stim} serves runs that include native noise.

\begin{algorithm}[htbp]
\caption{Randomness audit of a seeded QEC encoder under the bounded cloud fault model.}
\label{alg:qec:audit}
\KwIn{Code parameters $(n,k{=}1)$; HSE depth $d$; number of seeds $M$; fault set $\mathcal{F}$ (all $3n$ weight-one Paulis); adversary mode $\in\{\textsc{adaptive},\textsc{committed}\}$; validated DRBG}
\KwOut{Mean accepted logical disturbance $\bar\Delta$; fraction $\phi_{\mathrm{KL}}$ of encoders satisfying \cref{eq:qec:kl}}
$\Sigma\leftarrow 0$; $\Phi\leftarrow 0$\;
\lIf{mode $=$ \textsc{committed}}{adversary commits to $P^{\star}\in\mathcal{F}$ before any seed is drawn}
\For{$m\leftarrow 1$ \KwTo $M$}{
  $s_m\leftarrow\textsc{DRBG.Generate}()$ \tcp*{SP 800-90A instance inside a FIPS 140-3 module}
  $E_m\leftarrow\textsc{HSE}(s_m,d)$; derive $\mathcal{S}_m$, $\bar X_m$, $\bar Z_m$ by tableau conjugation \tcp*{\cref{eq:qec:encoder,eq:qec:hse}}
  encode $\ket{\psi}\ket{0^{\,n-1}}\mapsto E_m\big(\ket{\psi}\ket{0^{\,n-1}}\big)$ on the QPU\;
  \eIf{mode $=$ \textsc{adaptive}}{
    adversary learns $E_m$ and injects $P_m\leftarrow\argmax_{P\in\mathcal{F}}\Delta(E_m,P)$\;
  }{
    adversary injects $P_m\leftarrow P^{\star}$\;
  }
  measure the syndrome $\sigma(P_m)$ against $\mathcal{S}_m$; \textbf{accept} iff $\sigma=0$, otherwise \textbf{reject}\;
  $\Sigma\leftarrow\Sigma+\Delta(E_m,P_m)$ \tcp*{\cref{eq:qec:disturbance}; zero on rejection}
  $\Phi\leftarrow\Phi+\mathbb{1}\big[E_m\text{ satisfies \cref{eq:qec:kl} for }\mathcal{F}\big]$\;
  append $(m,\ \mathrm{id}(s_m),\ \sigma(P_m),\ \Delta(E_m,P_m))$ to the signed audit log\;
}
\Return $\bar\Delta\leftarrow\Sigma/M$, $\phi_{\mathrm{KL}}\leftarrow\Phi/M$\;
\end{algorithm}

\begin{table}[htbp]
\caption{Steps of the audit protocol and the randomness and assurance standards each step maps onto. The correspondence follows the standards context of \cite{guo2026qecaudit} and is developed further in \cref{ch:pqc}.}
\label{tab:qec:protocol}
\centering\footnotesize
\begin{tabularx}{\textwidth}{@{}c p{2.9cm} X p{3.9cm}@{}}
\toprule
Step & Operation & Requirement audited & Standard \\
\midrule
1 & Seed generation & Approved deterministic random-bit generator with a validated entropy source, instantiated inside a validated module & SP 800-90A~\cite{nist2015sp80090a}, SP 800-90B~\cite{nist2018sp80090b}, FIPS 140-3~\cite{nist2019fips1403} \\
2 & Seed transport to the cloud controller & Confidentiality of the seed against a quantum-capable eavesdropper & ML-KEM, FIPS 203~\cite{nist2024fips203} \\
3 & Encoder derivation $E_s=\textsc{HSE}(s,d)$ & Determinism and reproducibility from the seed; polynomial cost & \cref{sec:qec:hse} \\
4 & Encoding on the QPU & Start of the fault window & --- \\
5 & Fault injection (model) & Adaptive vs.\ committed adversary of \cref{ass:qec:model} & --- \\
6 & Syndrome extraction and rejection & Trusted measurement of the generators; accept iff $\sigma=0$ & --- \\
7 & Disturbance accounting & Exact enumeration (tableau) or sampled gate-level simulation (Stim) of $\Delta$ & \cref{eq:qec:disturbance} \\
8 & Attestation of the audit log & Integrity and non-repudiation of seeds, syndromes and disturbances & ML-DSA / SLH-DSA, FIPS 204/205~\cite{nist2024fips204,nist2024fips205} \\
\bottomrule
\end{tabularx}
\end{table}

The protocol is \emph{reproducible by construction}. Encoders are derived deterministically from logged seeds, so an auditor can regenerate every encoder, recompute every stabilizer and verify every reported syndrome. The artifact of \cite{guo2026qecaudit} is containerized and seeded for exactly this reason.

The protocol is also \emph{standards-facing} (\cref{sec:bg:crypto} reviews the relevant standards). The quality of the defense rests on the unpredictability of the seed, so the seed source must itself be an approved generator in a validated module. That means an SP 800-90A DRBG~\cite{nist2015sp80090a} fed by an SP 800-90B entropy source~\cite{nist2018sp80090b} and operated under FIPS 140-3~\cite{nist2019fips1403}. The seed's transport and the log's attestation are natural applications of the post-quantum standards ML-KEM and ML-DSA/SLH-DSA~\cite{nist2024fips203,nist2024fips204,nist2024fips205}.

The seed-derivation and encoder rows (steps 1 and 3) follow the standards discussion of \cite{guo2026qecaudit}. That discussion derives the layer choices from an SP 800-90A generator and takes FIPS 140-3 as the precedent for a validated module boundary. The transport and attestation rows (steps 2 and 8) are this dissertation's extension. \Cref{ch:pqc} describes the FIPS 140-3 module that would implement steps 1, 2 and 8.

\section{Experiments}\label{sec:qec:experiments}

\subsection{Setup}\label{sec:qec:setup}

Two experimental regimes were used in \cite{guo2026qecaudit}. In the \emph{exact} regime, encoders are Haar isometries, uniformly random Clifford circuits or HSE draws on $n=4$ to $8$ qubits with $k=1$. The fault set is all $3n$ weight-one Paulis, and $\Delta$ is evaluated exactly by dense linear algebra and \cref{lem:qec:indicator}. Native noise is absent, so every non-zero disturbance is attributable to the encoder ensemble and the adversary. The \code{5}{1}{3} code of \cref{tab:qec:code513} is the fixed control.

The adaptive adversary of \cref{ass:qec:model} (seed-aware in the paper's terms) evaluates all $3n$ faults against each encoder and injects the worst. The committed adversary (seed-blind) selects one fault on training seeds and is scored on disjoint held-out test seeds. The primary comparisons use $n=8$, $24$ faults, $200$ seeds per condition and HSE depths $d\in\{1,2,4,6,8,12\}$. The reported depth, $d=12$, maximizes the exact-correction fraction.

In the \emph{gate-level} regime, HSE encoders with $n\in\{5,9,13,17\}$ and $d\in\{2,4,8,12\}$ are compiled to gates and simulated with Stim~\cite{gidney2021stim}. Each condition uses $20$ seeds, every single-qubit Pauli fault and both logical input states. Ideal runs use one shot per circuit as a deterministic cross-check. Noisy runs use 1,000 shots under depolarizing noise with single-qubit rate $p_1=10^{-3}$ and two-qubit rate $p_2=10^{-2}$ applied after every encoder and decoder gate. Seeds were drawn from a DRBG in both regimes. Uncertainties are 95\% percentile-bootstrap intervals over seeds (5,000 resamples) and Wilson intervals for the exact-correction fractions.

\subsection{Results in the Exact Regime}\label{sec:qec:results513}

\Cref{tab:qec:results} summarizes the exact-regime results at the selected condition ($n=8$, $k=1$, $d=12$), and \cref{fig:qec:disturbance} plots them. Against the adaptive adversary, that is, a fixed and therefore learnable encoder, the mean accepted logical disturbance is $0.150$. By \cref{lem:qec:indicator}, $15\%$ of HSE draws have distance one, and the adversary finds the undetectable fault. Against the committed adversary the mean is $0.020$. For a fault chosen on training seeds and scored on held-out seeds, only $2\%$ of (encoder, fault) pairs let the fault through as a harmful logical.

The reduction is $(0.150-0.020)/0.150=86.7\%$, and it is attributable to \emph{rejection} (\cref{fig:qec:adversaries}). Under reseeding, a fault that is an undetectable logical for one encoder anticommutes with a stabilizer of the next and is caught. The 2-design reference of \cref{prop:qec:haar}, $0.0059$ for $n=8$, lies below both measured values, as it must. It is the average over all faults and encoders, whereas the audited numbers are a worst case over faults (adaptive) or a trained choice (committed). The measured Haar and uniform-Clifford means track the closed form to within $1.6\%$ at the largest width \cite{guo2026qecaudit}.

Finally, $18.5\%$ of the sampled depth-$12$ encoders satisfy the exact conditions \cref{eq:qec:kl} for the full weight-one fault set, against $14.5\%$ of uniformly random Clifford encoders. The uniform Clifford encoders use fewer two-qubit instructions at this width (median $33$ against $112$) but reach a lower exact-correction fraction. So nearly one draw in five is a genuine distance-three code, while the fixed \code{5}{1}{3} control corrects every tested fault with zero accepted disturbance.

\begin{table}[htbp]
\caption{Exact-regime audit results for eight-qubit HSE encoders of depth $12$ with all $24$ weight-one Pauli faults ($200$ seeds). Measured values are as reported in \cite{guo2026qecaudit}; the 2-design reference is computed from \cref{eq:qec:haar}.}
\label{tab:qec:results}
\centering\footnotesize
\begin{tabularx}{\textwidth}{@{}X c X@{}}
\toprule
Quantity & Value & Interpretation (\cref{lem:qec:indicator}) \\
\midrule
Encoder; fault set & HSE, $n=8$, $k=1$, $d=12$; $3n=24$ Paulis & all weight-one faults enumerated; \code{5}{1}{3} as fixed control \\
$\bar\Delta_{\mathrm{adapt}}$ (fixed encoder, seed-aware adversary) & $0.150$ & fraction of draws with a distance-one code \\
$\bar\Delta_{\mathrm{commit}}$ (reseeded encoder, seed-blind fault) & $0.020$ & fraction of (draw, fault) pairs accepted and harmful \\
Relative reduction & $86.7\%$ & attributable to syndrome-based rejection \\
Encoders satisfying \cref{eq:qec:kl} for $\mathcal{F}$ & $18.5\%$ ($14.5\%$ uniform Clifford) & distance-three codes among draws \\
2-design (Haar) reference, $n=8$ & $0.0059$ & $3\cdot 2^{n}/(2(4^{n}-1))$, \cref{prop:qec:haar} \\
Gate-level regime & $n\in\{5,9,13,17\}$; $p_1=10^{-3}$, $p_2=10^{-2}$ & Stim, 20 seeds per condition; scaling confirmed \\
\bottomrule
\end{tabularx}
\end{table}

% alt: Two-panel figure. Left, a bar chart with three bars: fixed encoder with seed-aware adversary at 0.150, reseeded encoder with seed-blind fault at 0.020, and the 2-design reference at 0.006. Right, a line chart of mean accepted logical disturbance versus HSE depth from 0 to 12 with two curves, adaptive and committed, both decreasing from their depth-zero values of 1.0 and 0.125 toward the reported values 0.150 and 0.020 at depth 12, with a dashed horizontal 2-design reference line at 0.006.
\begin{figure}[htbp]
\centering
\begin{subfigure}[t]{0.48\textwidth}
\centering
\begin{tikzpicture}
\begin{axis}[
  width=\linewidth,height=0.82\linewidth,
  ybar,bar width=7mm,
  ymin=0,ymax=0.2,ytick={0,0.05,0.10,0.15,0.20},
  ylabel={mean disturbance $\bar\Delta$},
  symbolic x coords={fixed,reseeded,2-design},
  xtick={fixed,reseeded,2-design},
  xticklabels={fixed\\adaptive,reseeded\\committed,2-design\\reference},
  xticklabel style={align=center,font=\footnotesize},
  yticklabel style={/pgf/number format/fixed},
  nodes near coords={\pgfmathprintnumber[fixed,precision=3]{\pgfplotspointmeta}},
  nodes near coords style={font=\footnotesize},
  enlarge x limits=0.4,
  cycle list={{fill=cbRed!70,draw=cbRed},{fill=cbBlue!70,draw=cbBlue},{fill=cbGray!35,draw=cbGray,densely dotted}},
  ]
\addplot coordinates {(fixed,0.150)};
\addplot coordinates {(reseeded,0.020)};
\addplot coordinates {(2-design,0.0059)};
\end{axis}
\end{tikzpicture}
\caption{Reported values for eight-qubit encoders ($d=12$).}
\label{fig:qec:disturbance-bars}
\end{subfigure}\hfill
\begin{subfigure}[t]{0.48\textwidth}
\centering
\begin{tikzpicture}
\begin{semilogyaxis}[
  width=\linewidth,height=0.82\linewidth,
  xlabel={HSE depth $d$ (layers)},
  ylabel={$\bar\Delta$},
  xmin=-0.5,xmax=12.5,ymin=0.004,ymax=1.5,
  xtick={0,2,4,6,8,10,12},
  legend style={at={(0.5,-0.30)},anchor=north,legend columns=1,font=\footnotesize,draw=none},legend cell align=left]
\addplot[cbRed,mark=diamond*,dashdotted,thick] coordinates {(0,1.0) (2,0.60) (4,0.38) (6,0.27) (8,0.21) (10,0.17) (12,0.150)};
\addlegendentry{adaptive (fixed)}
\addplot[cbBlue,mark=*,solid,thick] coordinates {(0,0.125) (2,0.075) (4,0.045) (6,0.032) (8,0.026) (10,0.022) (12,0.020)};
\addlegendentry{committed (reseeded)}
\addplot[cbGray,dashed,no marks,thick] coordinates {(-0.5,0.0059) (12.5,0.0059)};
\addlegendentry{2-design reference}
\end{semilogyaxis}
\end{tikzpicture}
\caption{Dependence on HSE depth (illustrative).}
\label{fig:qec:disturbance-depth}
\end{subfigure}
\caption{Accepted logical disturbance for eight-qubit HSE encoders with all $24$ weight-one faults. (a) Values reported in \cite{guo2026qecaudit} at depth $12$ for a fixed encoder facing a seed-aware adversary ($0.150$) and a reseeded encoder facing a seed-blind fault ($0.020$), with the 2-design reference $0.0059$ from \cref{eq:qec:haar}. (b) The depth-$0$ values ($1.0$ and $3/24=0.125$) are exact for the identity encoder and the depth-$12$ values are the reported ones; intermediate points are illustrative and are not measured data.}
\label{fig:qec:disturbance}
\end{figure}

The depth dependence in \cref{fig:qec:disturbance-depth} makes the mechanism visible. At depth zero the encoder is the identity, the ancillas are unprotected, and any fault on the logical qubit is accepted and harmful. The adaptive adversary then always wins ($\bar\Delta_{\mathrm{adapt}}=1$), and a committed fault is harmful in $3$ of $24$ cases ($\bar\Delta_{\mathrm{commit}}=0.125$), both exactly. Each additional layer spreads the logical operators over more qubits and converts more weight-one faults into detectable ones. The committed curve falls quickly, because a randomly placed fault rarely aligns with the few remaining weight-one logicals. The adaptive curve falls only as fast as the fraction of distance-one draws, because the adaptive adversary needs just one such logical to exist.

Because $\Delta$ is an indicator in this regime (\cref{lem:qec:indicator}), the audited means are Bernoulli proportions with elementary precision. With $M$ seeds the standard error of an estimated mean $\bar\Delta$ is $\sqrt{\bar\Delta(1-\bar\Delta)/M}$. At the $M=200$ seeds of \cite{guo2026qecaudit} this is $\pm0.010$ for $\bar\Delta_{\mathrm{commit}}=0.020$ and $\pm0.025$ for $\bar\Delta_{\mathrm{adapt}}=0.150$, which is why the paper reports 95\% bootstrap intervals. With $M=10^{4}$ seeds the two means would be resolved to $\pm0.0014$ and $\pm0.0036$. The audit is therefore practical: each seed costs one tableau computation of $\bigO{d\,n^{2}}$ operations and $3n$ syndrome checks. An operator can certify an ensemble to three significant figures in seconds on a laptop, and can re-certify it whenever the depth, the coupling graph or the fault set changes.

\subsection{Gate-Level Simulations at Seventeen Qubits}\label{sec:qec:results17}

The exact regime isolates the encoder but ignores noise. In the gate-level regime of \cite{guo2026qecaudit}, HSE encoders on $n=17$ qubits were compiled to gates and simulated in Stim with depolarizing noise of $p_1=10^{-3}$ and $p_2=10^{-2}$. Syndrome extraction was itself noisy, and the disturbance was estimated by sampling. \Cref{fig:qec:stim} sketches the qualitative picture, which is the one the exact regime predicts. Noise alone produces a logical disturbance that scales as a power of the physical rate set by the code distance of the draw. An adaptive adversary adds a floor that noise cannot lower, because the floor is the fraction of draws with an undetectable weight-one logical. Reseeding lowers that floor by the same mechanism as in the exact regime.

In the ideal gate-level runs, the fraction of single-qubit Pauli faults that cause an accepted logical change falls with both width and depth. It falls from about $0.15$ at $n=5$, $d=2$ to roughly $10^{-2}$ or less at $n=17$, $d=12$ (values read from Fig.~4 of \cite{guo2026qecaudit}). The ideal Stim classification agrees with the dense audit in $100\%$ of the overlapping cases. Under noise, the fraction of shots that are both accepted and corrupted falls the same way, from about $0.08$ at $n=5$, $d=2$ to the $10^{-6}$ range at $n=17$, $d=12$. There the mean acceptance is only $1.23\times10^{-5}$ and the unconditional accepted-failure rate $5.88\times10^{-6}$. The paper notes that this low failure rate records an availability collapse rather than a free lunch. At the widest and deepest condition almost every noisy shot is rejected, so a practical device must also meet an acceptance target over repeated fault-tolerant cycles.

Noisy syndrome extraction adds two effects that the exact regime cannot show. It can misreport a trivial syndrome as non-trivial, which raises the rejection rate and therefore lowers the accepted disturbance at the price of throughput. It can also misreport a non-trivial syndrome as trivial, which lets a fault through that the ideal check would have caught. Both effects grow with $p_2$ and with the depth of the extraction circuit, and both favour shallow encoders.

% alt: Log-log schematic of accepted logical disturbance versus two-qubit depolarizing rate for a seventeen-qubit register. A noise-only curve rises quadratically. A fixed-encoder-with-adaptive-adversary curve sits on a high floor and rises at large noise. A reseeded-encoder-with-committed-fault curve sits on a lower floor. A vertical dashed line marks the operating point p2 equals one percent.
\begin{figure}[htbp]
\centering
\begin{tikzpicture}
\begin{loglogaxis}[
  xlabel={two-qubit depolarizing rate $p_2$ (with $p_1=p_2/10$)},
  ylabel={accepted logical disturbance (schematic)},
  xmin=8e-4,xmax=7e-2,ymin=1e-5,ymax=20,
  legend pos=north west,legend cell align=left]
\addplot[cbRed,mark=diamond*,dashdotted,thick] coordinates {(1e-3,0.1000) (2e-3,0.1001) (5e-3,0.1008) (1e-2,0.103) (2e-2,0.112) (5e-2,0.175)};
\addlegendentry{fixed encoder, adaptive adversary}
\addplot[cbBlue,mark=*,solid,thick] coordinates {(1e-3,0.01003) (2e-3,0.01012) (5e-3,0.01075) (1e-2,0.013) (2e-2,0.022) (5e-2,0.085)};
\addlegendentry{reseeded encoder, committed fault}
\addplot[cbGreen,mark=triangle*,dotted,thick] coordinates {(1e-3,3e-5) (2e-3,1.2e-4) (5e-3,7.5e-4) (1e-2,3e-3) (2e-2,1.2e-2) (5e-2,7.5e-2)};
\addlegendentry{noise only, no adversary}
\addplot[cbGray,dashed,no marks] coordinates {(1e-2,1e-5) (1e-2,1)};
\node[anchor=south west,font=\footnotesize,cbGray] at (axis cs:1.05e-2,1.5e-5) {operating point $p_2=10^{-2}$};
\end{loglogaxis}
\end{tikzpicture}
\caption{Schematic of the gate-level regime studied at $n=17$ in \cite{guo2026qecaudit}: a noise-only disturbance that grows as a power of the physical error rate, an adversarial floor for a fixed encoder, and a lower floor after reseeding. The curves illustrate the qualitative form and are not measured values; only the operating point $p_1=10^{-3}$, $p_2=10^{-2}$ is taken from the paper.}
\label{fig:qec:stim}
\end{figure}

\section{Discussion and Limitations}\label{sec:qec:discussion}

\emph{Weight-one faults.} The bounded model allows one weight-one Pauli per run, the natural first case and the one a code is built for. But a co-tenant who can drive crosstalk on a coupler can plausibly produce correlated weight-two faults, and an adversary who controls timing can inject in more than one window. A weight-two fault set leaves \cref{lem:qec:indicator} and \cref{alg:qec:audit} unchanged in form, because the audit enumerates $\mathcal{F}$, whatever it is. But it raises the cost from $3n$ to $\bigO{n^{2}}$ faults per encoder and lowers the fraction of HSE draws that detect everything, so the disturbance values would rise. The Haar expectation of \cref{prop:qec:haar} holds for any traceless Hermitian unitary fault, so it remains the reference.

\emph{Small codes.} The exact experiments use $n\le 8$, where everything can be enumerated, and the gate-level experiments reach $n=17$. Neither is a size at which fault-tolerant computation will run. The method scales: tableau simulation is polynomial, Stim handles thousands of qubits, and the disturbance is defined for any code. The \emph{ensemble} does not obviously scale. At large $n$ a shallow HSE draw will rarely be a good code, so the operator will want to reseed within a family of encoders known to be good. One route randomizes the encoding circuit of a fixed surface or bicycle code~\cite{fowler2012surfacecodes,bravyi2024bicycle} through its automorphisms and gauge freedoms instead of drawing an arbitrary Clifford. The audit applies unchanged to such constrained ensembles.

\emph{Adversary model.} The adaptive adversary of \cref{ass:qec:model} is strong in knowledge (it learns the encoder exactly) and weak in capability (one Pauli). Real adversaries are the reverse: they see side channels rather than circuits, and they may cause more than one fault. The committed adversary, by contrast, is a faithful model of an attacker facing a fresh seed, provided the seed is unpredictable. This is why the standards mapping of \cref{tab:qec:protocol} is part of the defense rather than decoration. The model also treats syndrome extraction as trusted. An adversary who faults the extraction ancillas attacks a different surface, for which flag qubits and repeated extraction are the standard remedies.

\emph{Rejection versus correction.} Rejection is the cleanest primitive to audit, but computation needs correction. With correction, the quantity that survives is the disturbance after the decoder has applied its recovery. An adaptive adversary can then target faults whose syndrome the decoder maps to the wrong recovery. \Cref{def:qec:disturbance} extends directly: replace $\Pi_E$ with the decoder's recovery map. Reseeding helps for the same reason: the map from fault to syndrome changes with the seed. Quantifying this extension is the most important next step.

\emph{Throughput cost of rejection.} Rejection buys safety with discarded runs. In the exact regime the only rejected runs are those in which a fault was injected and caught, so the throughput cost is the adversary's injection rate. Under native noise, however, every run has a non-zero probability of a non-trivial syndrome. The acceptance rate then falls roughly as $(1-p_2)^{G_2}$, with $G_2$ the number of two-qubit gates in the encoder and extraction circuits. At $p_2=10^{-2}$ and a few dozen two-qubit gates this is a tens-of-percent overhead, tolerable for an audit but not for a long computation. The operational regime is correction, with rejection reserved for the audit itself and for the highest-assurance runs.

\emph{Cost of reseeding.} Every reseed costs a new compilation of the encoder and, on a fixed coupling graph, possibly a new routing. The HSE keeps this cheap because its layer structure is fixed and only the subsets, permutation and matching change. The learned synthesis of \cref{ch:synthesis} suggests a complementary route: a policy trained with a Clifford-fraction and topology-compliance reward emits compliant encoders for each seed directly.

\section{Relation to the Other Chapters}\label{sec:qec:relation}

This chapter attaches the dissertation's security thread to its error-correction thread and connects backward and forward. Backward, the Clifford-fraction and T-count terms of the \rubriq{} rubric in \cref{ch:synthesis} measure exactly the non-Clifford content whose cost the codes of this chapter impose on logic. The stabilizer tableau that \cref{alg:qec:audit} derives is the natural certificate of a synthesized encoder. A synthesizer that respects both constraints is the natural compiler for seeded, topology-compliant encoders. The gate-level simulations relied on Stim rather than on the GPU state-vector simulation of \cref{ch:qgear}, because stabilizer circuits admit polynomial simulation. The two simulators are complementary, and the non-Clifford logical operations of a full fault-tolerant computation would need the state-vector path.

Forward, the protocol of \cref{sec:qec:protocol} places three requirements on classical cryptography:
\begin{itemize}
  \item an approved DRBG in a validated module for the seed;
  \item a post-quantum key-encapsulation mechanism for its transport;
  \item a post-quantum signature for the audit log.
\end{itemize}
\Cref{ch:pqc} presents the FIPS 140-3 module with ML-KEM, ML-DSA and SLH-DSA support~\cite{guo2026lightriderqpc} that would meet them. The two chapters together answer RQ5. Error-corrected computation on shared hardware becomes auditable when its randomness is defined, measured and sourced to the same standard as the randomness in a cryptographic module.

\section{Summary}\label{sec:qec:summary}

We studied a security threat specific to error-corrected computation on shared cloud QPUs: a fixed, public encoder gives a co-located fault injector a reusable target. We formalized a bounded cloud fault model and defined the accepted logical disturbance, the acceptance-weighted harmful logical action that survives syndrome-based rejection. For Clifford encoders and Pauli faults it is an indicator of undetectable logicals, and its exact Haar expectation, $3\cdot 2^{n}/(2(4^{n}-1))$, holds for any unitary 2-design. We introduced the Hadamard-structured ensemble as a polynomial-cost seeded Clifford substitute for Haar encoders. We specified an audit protocol that is reproducible from logged seeds and maps onto SP 800-90A, FIPS 140-3 and FIPS 203/204/205.

For eight-qubit HSE encoders of depth $12$ with all $24$ weight-one faults, reseeding reduced the mean accepted logical disturbance from $0.150$ to $0.020$, an $86.7\%$ reduction attributable to rejection. Also, $18.5\%$ of the draws satisfied the exact error-correction conditions, the fixed \code{5}{1}{3} code corrected every tested fault, and gate-level Stim simulations up to $n=17$ under depolarizing noise confirmed the scaling~\cite{guo2026qecaudit}. The seed source is the hinge on which the whole defense turns, and the next chapter builds it.
